\section{Block-head candidate exclusion and the remaining bridge}

This chapter consolidates the passage from the block-head candidate
to the final valuation inequality.  It collects the exact identities
that connect the segment exclusions, the CRT candidate, and the
terminal residue.

\subsection{The block-head candidate}

The no-$H_{\ge}$ premises provide a block-head state of the form
\begin{equation*}
  r=\frac{A_j+5^jq}{2^{\wordWeight_j}}.
\end{equation*}
The unified core parameterises its predecessor by
\begin{equation*}
  x=2^{e-1}g+\delta5^{j-1},
  \qquad g=\Orbit_7(j-1),
\end{equation*}
with $e\ge2$ and $\delta\in\{1,3\}$.

\subsection{The q0 interval}

The block-head numerator satisfies the exact identity
\begin{equation*}
  A_j+5^jq=2^L\left(B+\delta5^j\right),
\qquad 0<B<5^j,\quad A_j<5^j.
\end{equation*}
Since $A_j<5^j$, the quotient structure forces
\begin{equation*}
  q=m+\delta\,2^L,
  \qquad m<2^L.
\end{equation*}
This is \texttt{reset\_\allowbreak q0\_\allowbreak form}.  The interval
condition
\begin{equation*}
  2^{W_p}\le q<2^{W_j}
\end{equation*}
is equivalent to the block-head bounds
\begin{equation*}
  5^j\le2^t r_j,\qquad r_j<5^j,
\end{equation*}
by \texttt{q0\_\allowbreak interval\_\allowbreak iff\_\allowbreak
rj\_\allowbreak bound}.

\subsection{The block-head identity}

Under the reset equation, the same numerator has the block-head form
\begin{equation*}
  A_j+5^jq
    =2^{W_p}\left(5^{k_0+1}s-4+\delta5^j\right).
\end{equation*}
This is \texttt{block\_\allowbreak head\_\allowbreak identity\_\allowbreak
of\_\allowbreak reset}.  It is the exact algebra that connects the
block ledger to the reset equation.

\subsection{Candidate emptiness}

The segment-length exclusions show that no candidate can satisfy the
full constraints:
\begin{enumerate}[leftmargin=*]
\item $d=1$ is excluded by the size bound $g<5^{j-1}/4$;
\item $d=2$ is excluded by the two congruences modulo $16$;
\item $d=3$ is excluded by the first-big-step contradiction;
\item $d\ge4$ is excluded by the finite base and the $e=3,j=17$
  residues.
\end{enumerate}
Hence the candidate family is empty.  The corresponding Lean
exclusions are compiled.

\subsection{From emptiness to the CRT candidate}

The exclusion layer leaves the CRT candidate $r_{j,0}$ and the
terminal residue $y^*$.  Let
\begin{equation*}
  n=s-j,\qquad \Delta=W_s-W_j,
\end{equation*}
and let $B$ be the block molecule of the tail of length $L$.  The tail
equation is
\begin{equation*}
  2^{\Delta}\,r_s=5^n\,r_{j,0}+B.
\end{equation*}
The terminal odd part is
\begin{equation*}
  w_{\mathrm{term}}=\frac{3r_s+1}{2^{L+4}},
\end{equation*}
and the terminal identity is
\begin{equation*}
  2^{\Delta+L+4}\,w_{\mathrm{term}}
    =3\cdot5^n\,r_{j,0}+3B+2^{\Delta}.
\end{equation*}

\subsection{The terminal valuation}

The final valuation inequality is
\begin{equation*}
  \vtwo{5^{L+3}w_{\mathrm{term}}+1}\le H_s-2.
\end{equation*}
It fails exactly when
\begin{equation*}
  w_{\mathrm{term}}\equiv-5^{-(L+3)}
  \pmod{2^{H_s-1}}.
\end{equation*}
Substituting this failure congruence into the terminal identity gives
the CRT equation
\begin{equation*}
  3\cdot5^n\,r_{j,0}+3B+2^{\Delta}
    =2^{\Delta+L+4}\left(u+m'\,2^{H_s-1}\right),
\end{equation*}
where $u$ is the least nonnegative residue of
$-5^{-(L+3)}$ modulo $2^{H_s-1}$ and $m'\ge0$.  The terminal residue
\begin{equation*}
  y^*=
  -\left(w_{\mathrm{term}}+5^{-(L+3)}\right)
  \left(3\cdot5^s\right)^{-1}
  \pmod{2^{H_s-1}}
\end{equation*}
is zero exactly in the failure case, so the inequality is equivalent
to $y^*\ne0$.

\subsection{The size conditions}

The two size conditions are
\begin{equation*}
  3B+2^{\Delta}\le3\cdot5^n-2\cdot4^n,
\end{equation*}
and
\begin{equation*}
  3\cdot5^s+\left(3\cdot5^n-2\cdot4^n\right)
    \le2^{\Delta+L+4}\,u.
\end{equation*}
They imply
\begin{equation*}
  r_{j,0}\ge\frac{3\cdot5^s}{3\cdot5^n}=5^j.
\end{equation*}
This is \texttt{rj0\_\allowbreak ge\_\allowbreak of\_\allowbreak
size\_\allowbreak conditions\_\allowbreak no\_\allowbreak hge},
compiled.

\subsection{The remaining bridge}

The lower bound $r_{j,0}\ge5^j$ must be compared with the interval
constraints carried by the no-$H_{\ge}$ premises.  The comparison is
the remaining Lean statement
\texttt{unified\_\allowbreak core\_\allowbreak final\_\allowbreak
no\_\allowbreak hge}.  It is the single \texttt{sorry} in the
unified-core audit; the paper does not claim it is closed.

\subsection{Statement inventory}

\begin{center}
\small
\begin{tabular}{@{}p{0.44\textwidth}p{0.28\textwidth}p{0.22\textwidth}@{}}
\toprule
Statement & File & Status \\
\midrule
\texttt{reset\_\allowbreak q0\_\allowbreak form} &
  \texttt{UnifiedCoreBridge.lean} & compiled \\
\texttt{q0\_\allowbreak interval\_\allowbreak iff\_\allowbreak
rj\_\allowbreak bound} &
  \texttt{UnifiedCoreBridge.lean} & compiled \\
\texttt{block\_\allowbreak head\_\allowbreak identity\_\allowbreak
of\_\allowbreak reset} &
  \texttt{UnifiedCoreBridge.lean} & compiled \\
segment exclusions &
  \texttt{UnifiedCoreBridge.lean} & compiled \\
\texttt{rj0\_\allowbreak spec\_\allowbreak eq} &
  \texttt{UnifiedCoreAudit.lean} & compiled \\
\texttt{terminal\_\allowbreak bound\_\allowbreak iff\_\allowbreak
yStar\_\allowbreak ne\_\allowbreak zero} &
  \texttt{UnifiedCoreAudit.lean} & compiled \\
\texttt{rj0\_\allowbreak ge\_\allowbreak of\_\allowbreak
size\_\allowbreak conditions\_\allowbreak no\_\allowbreak hge} &
  \texttt{UnifiedCoreAudit.lean} & compiled \\
\texttt{unified\_\allowbreak core\_\allowbreak final\_\allowbreak
no\_\allowbreak hge} &
  \texttt{UnifiedCoreAudit.lean} & pending \\
\bottomrule
\end{tabular}
\end{center}

\subsection{Status}

\begin{center}
\small
\begin{tabular}{@{}p{0.56\textwidth}p{0.36\textwidth}@{}}
\toprule
Layer & Status \\
\midrule
candidate parameterisation & math proven, Lean compiled \\
segment exclusions & math proven, Lean compiled \\
CRT candidate and terminal residue & math proven, Lean compiled \\
size-condition lower bound & math proven, Lean compiled \\
final valuation inequality & math proven, Lean pending \\
\bottomrule
\end{tabular}
\end{center}
